\chapter*{Abstract} \label{chap:abstract}

The increasing complexity and dynamic nature of modern robotic systems presents significant challenges regarding their analysis and Fault Detection. Manual model-based approaches are often too cost-, time-, and labour-intensive and thus do not match the concept of rapid prototyping and its swift, iterative development cycles. Moreover, the rapid changes in system design foster a rapid change in expert knowledge, which leads to inconsistency and thus degenerated system models, especially in systems that are maintained by multiple teams. This is compounded by the difficulty of creating dynamic system models that represent the actual runtime behaviour of a system, rather than simulations.

This thesis addresses these challenges by introducing a Digital Shadow --~a passive, data-driven digital representation of a system~-- that generates system specific models at runtime. To this end, multiple concepts of ROS2 are explored and abstracted by deriving a mathematical meta-model that enables the generation of dynamic system models that are capable of integrating platform-specific information. To enable faster feedback cycles and accelerate debugging in order to reduce development time and costs while increasing system reliability, the meta-model integrates several mechanisms to facilitate \acl{fdd}. These mechanisms aggregate the system to create holistic perspectives that provide a broad view of the system while keeping information fine-grained.

A Digital Shadow is employed in order to generate a digital representation of a system in question. It aims to reduce necessitated expert knowledge by employing a non-intrusive approach that leverages existing system data and runtime information. The Digital Shadow is designed to be lightweight in order to minimise its impact on the performance of the underlying system. Nevertheless, it captures both structural and continuous data to enable a comprehensive analysis and monitoring of the system's behaviour. To enable the subsequent analysis of the system, it additionally provides an interface for further applications that analyse the system in order to, e.g., perform model-based \acl{fdd}.

The approach is implemented in a proof-of-concept, mainly leveraging tracepoints in the ROS2 source code to achieve the objectives of this thesis. The feasibility of the mathematical meta-model is evaluated in a case study. Additionally, multiple experiments, which explore the performance of the proof-of-concept, are executed.

The results indicate that the proposed approach enables the dynamic and holistic modelling of ROS2 systems at runtime within defined constraints. However, within these constraints, the approach is capable of automatically generating ROS2 system specific models while minimising the required expert knowledge and negating the need for manipulating the software stack of the respective ROS2 system. Thereby, the overhead that is typically associated with manual modelling is reduced. Moreover, the Digital Shadow supports online \acl{fdd} by facilitating system analysis, aiding the development of more resilient, reliable, and maintainable robotic applications.